\documentclass[11pt,a4paper]{article}

% --- Packages ---
\usepackage[utf8]{inputenc}
\usepackage{amsmath,amssymb,amsfonts}
\usepackage{physics} % Critical for physics notation (bra-ket, derivatives, expectation values, etc.)
\usepackage[margin=1in]{geometry}
\usepackage{enumitem}
\usepackage[colorlinks=true, linkcolor=blue, urlcolor=cyan]{hyperref}
\usepackage{booktabs}
\usepackage{microtype}

% --- Document Metadata ---
\title{Quantum Computational Physics \\ \large Problem Statements}
\author{Numerical Methods in Quantum Mechanics}
\date{\today}

\begin{document}

\maketitle

\tableofcontents
\newpage

% --- Problem 1 ---
\section{Problem 1: The Truncated Parabola in a Harmonic Oscillator}
\label{sec:problem1}

\subsection*{Problem Statement}
Consider a particle of mass $m$ subject to a one-dimensional quantum harmonic oscillator potential:
\begin{equation}
    V(x) = \frac{1}{2}m\omega^2x^2
\end{equation}

At time $t=0$, the particle is prepared in the initial state:
\begin{equation}
    \psi(x,0) = 
    \begin{cases} 
        N(A - Bx^2) & \text{for } x^2 \le \frac{A}{B} \\ 
        0 & \text{for } x^2 > \frac{A}{B} 
    \end{cases}
\end{equation}
where $A$ and $B$ are positive real constants, and $N$ is the normalization constant.

\subsection*{Analytical Tasks}
\begin{enumerate}[label=\textbf{\arabic*.}]
    \item \textbf{Normalization Constant:} Determine the normalization constant $N$ in terms of $A$ and $B$ by enforcing the normalization condition:
    \begin{equation}
        \int_{-\infty}^{\infty} |\psi(x,0)|^2 \dd{x} = 1
    \end{equation}
    
    \item \textbf{Eigenstate Expansion:} Calculate the probability $P_n = |c_n|^2$ of finding the particle in the energy eigenstate $\ket{\phi_n}$ with energy eigenvalue $E_n = \left(n + \frac{1}{2}\right)\hbar\omega$. The expansion coefficients are given by the projection:
    \begin{equation}
        c_n = \braket{\phi_n}{\psi(x,0)} = \int_{-\sqrt{A/B}}^{\sqrt{A/B}} \phi_n^*(x) \psi(x,0) \dd{x}
    \end{equation}
    where $\phi_n(x)$ are the Hermite-Gaussian eigenfunctions of the quantum harmonic oscillator.

    \item \textbf{Time Evolution:} Write down the formal expression for the time-evolved wave function $\psi(x,t)$ for $t > 0$:
    \begin{equation}
        \psi(x,t) = \sum_{n=0}^{\infty} c_n \phi_n(x) e^{-i E_n t / \hbar}
    \end{equation}
\end{enumerate}

\newpage

% --- Problem 2 ---
\section{Problem 2: The Shifted and Squeezed Gaussian}
\label{sec:problem2}

\subsection*{Problem Statement}
A particle of mass $m$ in a harmonic oscillator potential $V(x) = \frac{1}{2}m\omega^2x^2$ is initially prepared in a displaced (shifted) and squeezed Gaussian wave packet:
\begin{equation}
    \psi(x,0) = N e^{-\frac{(x-x_0)^2}{a^2}}
\end{equation}
where:
\begin{itemize}
    \item $x_0$ is the initial spatial displacement (shift) of the wave packet center.
    \item $a$ is the initial width of the wave packet, such that $a \neq \sqrt{\frac{\hbar}{m\omega}}$ (meaning the wave packet width differs from the ground-state characteristic length scale $x_c = \sqrt{\frac{\hbar}{m\omega}}$, resulting in a squeezed state).
    \item $N$ is the normalization constant.
\end{itemize}

\subsection*{Analytical Tasks}
\begin{enumerate}[label=\textbf{\arabic*.}]
    \item \textbf{Time Evolution:} Determine the full time-evolved wave packet $\psi(x,t)$ under the harmonic oscillator Hamiltonian. Show that the wave packet remains Gaussian in form:
    \begin{equation}
        \psi(x,t) = N(t) e^{-A(t)(x - \langle x \rangle(t))^2 + \frac{i}{\hbar} \langle p \rangle(t) x + i \phi(t)}
    \end{equation}
    where $A(t)$ is a complex parameter related to the time-dependent width and phase.

    \item \textbf{Expectation Values:} Calculate the expectation values of position $\langle x \rangle(t)$ and momentum $\langle p \rangle(t)$ as functions of time. Verify that they satisfy Ehrenfest's theorem and display classical oscillatory behavior:
    \begin{align}
        \langle x \rangle(t) &= x_0 \cos(\omega t) \\
        \langle p \rangle(t) &= -m\omega x_0 \sin(\omega t)
    \end{align}

    \item \textbf{Uncertainty Relations and Breathing Mode:} 
    Determine the position uncertainty $\Delta x(t) = \sqrt{\langle x^2 \rangle - \langle x \rangle^2}$ and momentum uncertainty $\Delta p(t)$. Show that the wave packet width oscillates (or "breathes") at a frequency of $2\omega$.
\end{enumerate}

\newpage

% --- Problem 3 ---
\section{Problem 3: Tunneling in the Quartic Double Well}
\label{sec:problem3}

\subsection*{Problem Statement}
A particle of mass $m$ moves in a one-dimensional quartic double-well potential described by:
\begin{equation}
    V(x) = -\frac{1}{2}\alpha x^2 + \beta x^4
\end{equation}
where $\alpha$ and $\beta$ are positive constants. This potential features two symmetric minima at $x_{\text{min}} = \pm \sqrt{\frac{\alpha}{4\beta}}$ separated by a potential barrier of height $V_0 = \frac{\alpha^2}{16\beta}$ at $x = 0$.

\subsection*{Physical and Computational Tasks}
\begin{enumerate}[label=\textbf{\arabic*.}]
    \item \textbf{Eigenstate Calculation:} Find the lowest energy eigenstates by solving the time-independent Schrödinger equation:
    \begin{equation}
        \left[ -\frac{\hbar^2}{2m} \frac{\dd^2}{\dd{x}^2} + V(x) \right] \psi_n(x) = E_n \psi_n(x)
    \end{equation}
    numerically on a discretized spatial grid of size $N$ over the domain $x \in [-L, L]$ using finite difference approximations.

    \item \textbf{Energy Eigenvalues and Degeneracy Splitting:} 
    Determine the ground state energy $E_0$ (symmetric state, even parity) and the first excited state energy $E_1$ (antisymmetric state, odd parity). Calculate the energy splitting:
    \begin{equation}
        \Delta E = E_1 - E_0
    \end{equation}
    Explain how quantum tunneling lifts the classical degeneracy of the two isolated wells.

    \item \textbf{Tunneling Time:} 
    Construct a state initially localized in the right well:
    \begin{equation}
        \psi_{R}(x,0) = \frac{1}{\sqrt{2}} \left( \psi_0(x) + \psi_1(x) \right)
    \end{equation}
    Show that the wave function evolves in time as:
    \begin{equation}
        \psi(x,t) = \frac{1}{\sqrt{2}} \left( \psi_0(x)e^{-iE_0t/\hbar} + \psi_1(x)e^{-iE_1t/\hbar} \right)
    \end{equation}
    Calculate the tunneling time $\tau$ required for the particle to tunnel completely from the right well to the left well, where:
    \begin{equation}
        \tau = \frac{\pi\hbar}{\Delta E}
    \end{equation}

    \item \textbf{Grid Convergence Study:}
    Conduct a numerical grid convergence test by solving the eigenvalue problem for varying spatial discretization grid points $N \in \{200, 400, 800, 1600, 3200\}$ to verify the stability and accuracy of the computed eigenvalues $E_0$ and energy splitting $\Delta E$.
\end{enumerate}
\newpage
%--Problem 4----
\section{Problem 4: Numerical Solution of the Schrödinger Equation}
\label{sec:problem4}

This section focuses on the numerical solution of the Time-Independent Schrödinger Equation (TISE) for different potentials, followed by the analysis of quantum dynamics using the computed eigenstates.

\subsection{Problem 1: Quantum Harmonic Oscillator}

\begin{enumerate}
    \item Compute the first 20 eigenvalues (energy levels) and their corresponding eigenfunctions of the quantum harmonic oscillator.
    \item Compare the computed eigenvalues with the exact analytical solutions and comment on the accuracy of the numerical method.
\end{enumerate}

\subsection{Problem 2: Quartic Potential}

\begin{enumerate}
    \item Compute the eigenvalues and eigenfunctions of the quartic potential using the same numerical method.
    \item Verify the numerical consistency of the solution by varying the computational grid size and examining the convergence of the eigenvalues and eigenfunctions.
\end{enumerate}

\subsection{Problem 3: Gaussian Wave Packet Evolution in a Double-Well Potential}

\begin{enumerate}
    \item Consider a Gaussian wave packet initially localized near the left minimum of the double-well potential.
    \item Expand the Gaussian wave packet in terms of the first 20 energy eigenstates.
    \item Examine whether the expansion coefficients for higher-energy states become negligibly small as the quantum number increases. If not, include additional eigenstates until the expansion converges.
    \item Using the converged expansion, compute and plot the time evolution of the wavefunction, $\psi(x,t)$.
    \item Analyze the resulting dynamics and identify the quantum tunneling effect. Comment on how the tunneling behavior changes for a deeper double-well potential.
\end{enumerate}

\end{document}
